% ============================================================
%  CHAPTER 2: SCALARS AND VECTORS
%  The Physics Codex: A Complete University Foundation
%  Korede Samuel Omotosho (Falcon)
%  Aurikrex Learning Systems, 2026
%
%  File: chapters/part1/ch02_scalars_vectors.tex
% ============================================================

\chapter{Scalars and Vectors}
\label{ch:scalars_vectors}

\chapterquote{The laws of nature are but the mathematical
thoughts of God.}{Euclid}

\begin{objectivebox}
By the end of this chapter, you should be able to:
\begin{enumerate}[noitemsep]
  \item Distinguish clearly between scalar and vector
  quantities with examples
  \item Identify and classify the different types of vectors
  \item Represent vectors graphically using arrows
  \item Calculate the magnitude of a vector in 2D and 3D
  \item Add vectors using the triangle, parallelogram,
  and polygon rules
  \item Resolve a vector into perpendicular components
  \item Add multiple vectors by resolving into components
  \item Define and calculate the resultant of two or more
  vectors
  \item Subtract vectors and find relative velocity
  \item Express vectors using unit vectors and position vectors
  \item Calculate the scalar (dot) product and use it to
  find the angle between two vectors
  \item Calculate the vector (cross) product and apply
  the right-hand rule
  \item Apply vector methods to solve real physical problems
\end{enumerate}
\end{objectivebox}

\vspace{0.4cm}

\minitoc

\vspace{0.5cm}

% ============================================================
\section{Revisiting Scalars and Vectors}
% ============================================================

In Chapter~1, we introduced scalars and vectors briefly.
Now we treat them in full. This is one of the most important
chapters in the entire textbook because \textit{every}
subsequent topic --- forces, motion, energy, electricity,
magnetism --- requires you to handle vectors correctly.

A vector is not simply ``a number with a direction''
attached as an afterthought. The direction is as
fundamental as the magnitude. Two vectors with the same
magnitude but different directions are completely different
physical quantities --- as different as a car heading north
at 100 km\,h$^{-1}$ is from a car heading south at the
same speed.

\begin{definitionbox}[Scalar]
A \textbf{scalar} is a quantity fully described by a
magnitude and a unit alone. Scalars obey ordinary
arithmetic. Examples: mass ($m = 5.0\ \text{kg}$),
temperature ($T = 300\ \text{K}$), energy
($E = 200\ \text{J}$), speed ($v = 15\ \text{m\,s}^{-1}$),
distance ($d = 40\ \text{m}$), time ($t = 5\ \text{s}$),
pressure, density, work, power.
\end{definitionbox}

\begin{definitionbox}[Vector]
A \textbf{vector} is a quantity requiring both magnitude,
unit, and direction for complete description. Vectors obey
the laws of vector algebra. Examples: displacement
($\vect{s}$), velocity ($\vect{v}$), acceleration
($\vect{a}$), force ($\vect{F}$), weight ($\vect{W}$),
momentum ($\vect{p}$), electric field ($\vect{E}$),
magnetic flux density ($\vect{B}$).
\end{definitionbox}

\begin{keybox}[Notation]
In this textbook, vectors are written in bold:
$\vect{F}$, $\vect{v}$, $\vect{a}$.
The magnitude of vector $\vect{F}$ is written as $F$
or $|\vect{F}|$. In handwriting, draw an arrow above
the symbol: $\vec{F}$.
\end{keybox}

% ============================================================
\section{Types of Vectors}
\label{sec:types_of_vectors}
% ============================================================

Not all vectors are created equal. Physicists and
mathematicians have given names to vectors with
particular properties. Recognising these types will
help you communicate precisely and avoid errors when
solving problems.

% ----------------------------------------------------------
\subsection*{Equal Vectors}

\textbf{Intuition:} Two vectors are equal if they are
completely interchangeable --- same size, same direction.

\begin{definitionbox}[Equal Vectors]
Two vectors $\vect{A}$ and $\vect{B}$ are \textbf{equal}
if and only if they have the same magnitude \textit{and}
the same direction:
\[
  \vect{A} = \vect{B}
  \iff |\vect{A}| = |\vect{B}|
  \text{ and same direction.}
\]
Their position in space does not matter --- equal vectors
can be moved anywhere (see \textit{free vectors} below).
\end{definitionbox}

\begin{diagbox}[Equal Vectors]
\begin{center}
\begin{tikzpicture}[scale=0.85]
  \draw[->, thick, codexblue, line width=2pt]
    (0,0) -- (2.5,1) node[midway, below]{$\vect{A}$};
  \draw[->, thick, codexblue, line width=2pt]
    (3.5,0.5) -- (6,1.5) node[midway, below]{$\vect{B}$};
  \node at (3,2) {\small Same length, same direction
    $\Rightarrow \vect{A} = \vect{B}$};
\end{tikzpicture}
\end{center}
\end{diagbox}

% ----------------------------------------------------------
\subsection*{Null (Zero) Vector}

\textbf{Intuition:} A displacement of zero --- you move
nowhere. Its arrow has zero length.

\begin{definitionbox}[Null Vector]
The \textbf{null vector} $\vect{0}$ has magnitude zero.
Its direction is indeterminate (undefined).
\[
  \vect{A} + (-\vect{A}) = \vect{0}
\]
Physical meaning: zero displacement, zero force, zero
velocity.
\end{definitionbox}

% ----------------------------------------------------------
\subsection*{Negative Vector}

\textbf{Intuition:} Reversing a journey. Same road,
opposite end.

\begin{definitionbox}[Negative Vector]
The \textbf{negative} of vector $\vect{A}$, written
$-\vect{A}$, has the same magnitude as $\vect{A}$ but
points in the exactly opposite direction:
\[
  |-\vect{A}| = |\vect{A}|,\quad
  \text{direction of } {-\vect{A}}
  = \text{direction of } \vect{A} + 180°.
\]
\end{definitionbox}

% ----------------------------------------------------------
\subsection*{Parallel and Anti-parallel Vectors}

\begin{definitionbox}[Parallel Vectors]
In the broad geometric sense, parallel vectors are collinear:
they point in the same or exactly opposite directions. In this
chapter, ``parallel'' without the prefix ``anti-'' means
\textbf{same direction} (angle between them $= 0°$).
They may have different magnitudes.
\end{definitionbox}

\begin{definitionbox}[Anti-parallel Vectors]
Two vectors are \textbf{anti-parallel} if they point
in exactly opposite directions (angle between them
$= 180°$). A negative vector is always anti-parallel
to the original.
\end{definitionbox}

\begin{diagbox}[Parallel and Anti-parallel]
\begin{center}
\begin{tikzpicture}[scale=0.85]
  % Parallel
  \draw[->, thick, codexblue, line width=2pt]
    (0,1.5) -- (2.5,1.5) node[midway, above]{$\vect{A}$};
  \draw[->, thick, codexred, line width=2pt]
    (0,0.8) -- (1.5,0.8) node[midway, below]{$\vect{B}$};
  \node at (1.25,2.2) {\small Parallel};

  % Anti-parallel
  \draw[->, thick, codexblue, line width=2pt]
    (5,1.5) -- (7.5,1.5) node[midway, above]{$\vect{C}$};
  \draw[->, thick, codexred, line width=2pt]
    (7.5,0.8) -- (5,0.8) node[midway, below]{$\vect{D}$};
  \node at (6.25,2.2) {\small Anti-parallel};
\end{tikzpicture}
\end{center}
\end{diagbox}

% ----------------------------------------------------------
\subsection*{Unit Vector}

\begin{definitionbox}[Unit Vector]
A \textbf{unit vector} is a vector whose magnitude is
exactly 1. It carries direction information only ---
no size. The hat notation $\hat{A}$ always denotes a
unit vector in the direction of $\vect{A}$.
\end{definitionbox}

\noindent The standard Cartesian unit vectors are
$\hat{i}$, $\hat{j}$, $\hat{k}$ along the $x$, $y$,
$z$ axes respectively. Constructing a unit vector from
any given vector is treated fully in
Section~\ref{sec:unit_vectors}.

% ----------------------------------------------------------
\subsection*{Position Vector}

\begin{definitionbox}[Position Vector]
The \textbf{position vector} $\vect{r}$ of a point $P$
is the vector drawn from the origin $O$ to $P$.
It locates a point in space relative to a fixed
reference origin.
\[
  \vect{r} = x\hat{i} + y\hat{j}
  \quad \text{(2D)}
  \qquad
  \vect{r} = x\hat{i} + y\hat{j} + z\hat{k}
  \quad \text{(3D)}
\]
\end{definitionbox}

\noindent Position vectors are developed fully in
Section~\ref{sec:position_vectors}.

% ----------------------------------------------------------
\subsection*{Free Vectors}

\begin{definitionbox}[Free Vector]
A \textbf{free vector} is one that can be moved to any
position in space without changing its physical meaning,
provided its magnitude and direction are preserved.
\end{definitionbox}

\noindent Most vectors in introductory mechanics ---
forces (when only their net effect matters),
displacements, velocities --- are treated as free
vectors. This is why the triangle rule allows you to
slide vectors around to join them head-to-tail.

% ----------------------------------------------------------
\subsection*{Localised (Bound) Vector}

\begin{definitionbox}[Localised Vector]
A \textbf{localised} (or \textbf{bound}) vector has a
fixed point of application. Its effect depends on
\textit{where} it acts, not just its magnitude and
direction.
\end{definitionbox}

\noindent Example: The torque produced by a force
depends on where the force is applied relative to
the pivot. The same force applied at a different
point produces a different torque. Localised vectors
appear in rotational mechanics and structural analysis.

% ----------------------------------------------------------
\subsection*{Collinear Vectors}

\begin{definitionbox}[Collinear Vectors]
Vectors are \textbf{collinear} if they all lie along the
same straight line (or parallel lines). Parallel and
anti-parallel vectors are special cases of collinear
vectors.
\end{definitionbox}

% ----------------------------------------------------------
\subsection*{Coplanar Vectors}

\begin{definitionbox}[Coplanar Vectors]
Vectors are \textbf{coplanar} if they all lie in the
same plane. Any two vectors are automatically coplanar.
Three or more vectors may or may not be coplanar.
\end{definitionbox}

\noindent All vector problems in this chapter involve
coplanar vectors (2D). Three-dimensional, non-coplanar
vectors appear in later chapters on electromagnetism
and rotational dynamics.

\begin{keybox}[Summary of Vector Types]
\centering
\begin{tabular}{ll}
\toprule
\textbf{Type} & \textbf{Defining property}\\
\midrule
Equal & Same magnitude AND same direction\\
Null  & Magnitude $= 0$\\
Negative & Same magnitude, opposite direction\\
Parallel & Same direction (any magnitudes)\\
Anti-parallel & Opposite directions (any magnitudes)\\
Unit & Magnitude $= 1$ (direction only)\\
Position & Drawn from the origin to a point\\
Free & Can be translated anywhere\\
Localised & Fixed point of application\\
Collinear & Lie along the same line/direction\\
Coplanar & Lie in the same plane\\
\bottomrule
\end{tabular}
\end{keybox}

% ============================================================
\section{Magnitude of a Vector}
\label{sec:magnitude}
% ============================================================

\textbf{Intuition:} Magnitude is simply the \textit{size}
of a vector --- how long its arrow is. If you draw a
vector to scale, the magnitude is what you measure with
a ruler (in the appropriate units). Magnitude is always
a non-negative scalar.

% ----------------------------------------------------------
\subsection*{Derivation from Pythagoras's Theorem}

Consider a 2D vector $\vect{A} = A_x\hat{i} + A_y\hat{j}$.
Its components $A_x$ and $A_y$ form two sides of a
right-angled triangle, with the vector itself as the
hypotenuse. By Pythagoras:

\[
  |\vect{A}|^2 = A_x^2 + A_y^2
\]

\begin{diagbox}[Magnitude from Pythagoras]
\begin{center}
\begin{tikzpicture}[scale=1.0]
  \draw[->, thick, codexblue, line width=2.5pt]
    (0,0) -- (4,3)
    node[above right]{$\vect{A}$};
  \draw[->, thick, codexred, line width=2pt]
    (0,0) -- (4,0)
    node[below]{$A_x$};
  \draw[->, thick, codexgreen, line width=2pt]
    (4,0) -- (4,3)
    node[right]{$A_y$};
  \draw[gray, dashed] (0,3) -- (4,3);
  \draw[gray] (3.8,0) -- (3.8,0.2) -- (4,0.2);
  \draw[codexgold] (0.8,0) arc(0:36.87:0.8);
  \node[codexgold] at (1.15,0.28) {$\theta$};
  \node at (2,-1.0) {\small $|\vect{A}| = \sqrt{A_x^2 + A_y^2}$};
\end{tikzpicture}
\end{center}
\end{diagbox}

In three dimensions, extend Pythagoras once more along
the $z$-axis:

\begin{lawbox}[Magnitude of a Vector]
\textbf{2D:}\quad
$|\vect{A}| = \sqrt{A_x^2 + A_y^2}$

\vspace{0.3cm}

\textbf{3D:}\quad
$|\vect{A}| = \sqrt{A_x^2 + A_y^2 + A_z^2}$
\end{lawbox}

% ----------------------------------------------------------
\subsection*{Range of the Resultant}

A very useful theorem follows directly from the idea
of magnitude:

\begin{lawbox}[Range of the Resultant]
For two vectors of magnitudes $A$ and $B$, the
magnitude $R$ of their resultant must satisfy:
\[
  |A - B| \leq R \leq A + B
\]
\textbf{Maximum} $R = A + B$: when the vectors are
parallel ($\theta = 0°$).\\
\textbf{Minimum} $R = |A - B|$: when the vectors are
anti-parallel ($\theta = 180°$).\\
Any value of $R$ outside this range is physically
impossible.
\end{lawbox}

\noindent \textit{Physical intuition:} If one person
pulls a rope north with $10\ \text{N}$ and another
pulls east with $6\ \text{N}$, the resultant cannot
be $17\ \text{N}$ (that would require them to be
pulling in the same direction) and cannot be
$3\ \text{N}$ (that would require them to be pulling
against each other with a difference of only $3\ \text{N}$,
but $10 - 6 = 4\ \text{N}$ is the actual minimum).

\begin{examplebox}[2.1]
\stepcounter{workedexample}
Find the magnitude of the vector
$\vect{A} = 3\hat{i} + 4\hat{j}\ \text{N}$.
\end{examplebox}

\begin{solutionbox}
\[
  |\vect{A}| = \sqrt{A_x^2 + A_y^2}
  = \sqrt{3^2 + 4^2}
  = \sqrt{9 + 16}
  = \sqrt{25} = 5\ \text{N}
\]
\end{solutionbox}

\begin{examplebox}[2.2]
\stepcounter{workedexample}
A velocity vector in 3D is
$\vect{v} = 2\hat{i} - 3\hat{j} + 6\hat{k}\ \text{m\,s}^{-1}$.
Find its magnitude.
\end{examplebox}

\begin{solutionbox}
\[
  |\vect{v}| = \sqrt{2^2 + (-3)^2 + 6^2}
  = \sqrt{4 + 9 + 36}
  = \sqrt{49}
  = 7\ \text{m\,s}^{-1}
\]
\end{solutionbox}

% ============================================================
\section{Graphical Representation of Vectors}
% ============================================================

A vector is represented graphically by an arrow. The
\textbf{length} of the arrow represents the magnitude
(drawn to scale), and the \textbf{direction} of the arrow
represents the direction of the vector.

\begin{diagbox}[Vector Representation]
\begin{center}
\begin{tikzpicture}[scale=1.0]
  % Single vector
  \draw[->, thick, codexblue, line width=2pt]
    (0,0) -- (4,0);
  \node[below] at (2,0) {magnitude = $4\,\mathrm{N}$};
  \node[above] at (2,0.05) {$\vect{F}$};
  \draw[<->, gray] (0,-0.5) -- (4,-0.5);
  \node[below, gray] at (2,-0.5) {Scale: 1 cm = 1 N};

  % Direction angle
  \draw[->, thick, codexred, line width=2pt]
    (6,0) -- ++(3,2);
  \node at (7.8,0.9) {$\vect{v}$};
  \draw[gray] (6,0) -- ++(1.5,0);
  \draw[codexgold] (6.6,0) arc (0:33.7:0.6);
  \node[right, codexgold] at (7.0, 0.2) {$\theta$};
  \node[below] at (7.5,-0.2) {Direction: $\theta$ above horizontal};
\end{tikzpicture}
\end{center}
\end{diagbox}

% ============================================================
\section{Vector Addition: Triangle Rule}
% ============================================================

\begin{definitionbox}[Resultant Vector]
The \textbf{resultant} $\vect{R}$ of two or more vectors
is the single vector that has the same effect as all the
original vectors acting together. It is found by vector
addition.
\end{definitionbox}

\begin{lawbox}[Triangle Rule for Vector Addition]
To add two vectors $\vect{A}$ and $\vect{B}$:
\begin{enumerate}[noitemsep]
  \item Draw $\vect{A}$ as an arrow to scale.
  \item From the \textit{tip} (head) of $\vect{A}$, draw
  $\vect{B}$ to scale.
  \item The resultant $\vect{R} = \vect{A} + \vect{B}$
  is the arrow from the \textit{tail} of $\vect{A}$ to
  the \textit{head} of $\vect{B}$.
\end{enumerate}
This is called the \textbf{head-to-tail} method.
\end{lawbox}

\begin{diagbox}[Triangle Rule]
\begin{center}
\begin{tikzpicture}[scale=0.9]
  % Vector A
  \draw[->, thick, codexblue, line width=2pt]
    (0,0) -- (4,0) node[midway, below]{$\vect{A}$};
  % Vector B from tip of A
  \draw[->, thick, codexred, line width=2pt]
    (4,0) -- (6,2) node[midway, right]{$\vect{B}$};
  % Resultant
  \draw[->, thick, codexgreen, line width=2.5pt, dashed]
    (0,0) -- (6,2) node[midway, above left]{$\vect{R} = \vect{A}+\vect{B}$};
  % Labels
  \fill[codexblue] (0,0) circle (2pt) node[left]{Tail of A};
  \fill[codexred]  (6,2) circle (2pt) node[above right, xshift=2pt, yshift=2pt]{Head of B};
\end{tikzpicture}
\end{center}
\end{diagbox}

\textbf{Important properties of vector addition:}
\begin{align*}
  \vect{A} + \vect{B} &= \vect{B} + \vect{A}
  \quad \text{(commutative)}\\
  (\vect{A}+\vect{B})+\vect{C}
  &= \vect{A}+(\vect{B}+\vect{C})
  \quad \text{(associative)}
\end{align*}

% ============================================================
\section{Vector Addition: Parallelogram Rule}
% ============================================================

\begin{lawbox}[Parallelogram Rule]
To add vectors $\vect{A}$ and $\vect{B}$ using the
parallelogram method:
\begin{enumerate}[noitemsep]
  \item Draw both vectors from the same point (tail to tail).
  \item Complete the parallelogram.
  \item The diagonal of the parallelogram from the common
  tail is the resultant $\vect{R}$.
\end{enumerate}
\end{lawbox}

\begin{diagbox}[Parallelogram Rule]
\begin{center}
\begin{tikzpicture}[scale=0.9]
  \coordinate (O) at (0,0);
  \coordinate (A) at (4,0);
  \coordinate (B) at (1.5,2.5);
  \coordinate (R) at (5.5,2.5);

  % Parallelogram sides
  \draw[dashed, gray] (A) -- (R);
  \draw[dashed, gray] (B) -- (R);

  % Vectors
  \draw[->, thick, codexblue, line width=2pt]
    (O) -- (A) node[midway, below]{$\vect{A}$};
  \draw[->, thick, codexred, line width=2pt]
    (O) -- (B) node[midway, left]{$\vect{B}$};

  % Resultant diagonal
  \draw[->, thick, codexgreen, line width=2.5pt]
    (O) -- (R) node[midway, above, sloped]{$\vect{R}$};

  \fill (O) circle (2.5pt) node[left]{O};
\end{tikzpicture}
\end{center}
\end{diagbox}

% ============================================================
\section{Vector Addition: Polygon Rule}
\label{sec:polygon_rule}
% ============================================================

The triangle rule handles two vectors. What if you have
three, four, or more? Simply extend the head-to-tail
idea: keep chaining vectors end-to-end until all have
been placed, then draw the resultant from the tail of
the first to the head of the last.

\begin{lawbox}[Polygon Rule for Vector Addition]
To add $n$ vectors $\vect{F}_1, \vect{F}_2, \ldots,
\vect{F}_n$:
\begin{enumerate}[noitemsep]
  \item Draw $\vect{F}_1$ to scale.
  \item From the head of $\vect{F}_1$, draw $\vect{F}_2$.
  \item Continue placing each vector head-to-tail.
  \item The resultant $\vect{R}$ is the vector from the
  tail of $\vect{F}_1$ to the head of $\vect{F}_n$.
\end{enumerate}
\textbf{Equilibrium condition:} If the polygon
\textit{closes} (the head of the last vector meets the
tail of the first), then the resultant is zero and the
system is in equilibrium.
\end{lawbox}

\begin{diagbox}[Polygon Rule: Four Vectors]
\begin{center}
\begin{tikzpicture}[scale=0.85]
  \coordinate (O)  at (0,0);
  \coordinate (P1) at (3,0);
  \coordinate (P2) at (4.5,2);
  \coordinate (P3) at (2.5,3.5);

  % The four vectors head-to-tail
  \draw[->, thick, codexblue, line width=2pt]
    (O) -- (P1) node[midway, below]{$\vect{F}_1$};
  \draw[->, thick, codexred, line width=2pt]
    (P1) -- (P2) node[midway, right]{$\vect{F}_2$};
  \draw[->, thick, codexgold, line width=2pt]
    (P2) -- (P3) node[midway, right]{$\vect{F}_3$};
  \draw[->, thick, violet, line width=2pt]
    (P3) -- (1.0,4.5)
    node[midway, above right, xshift=2pt, yshift=2pt]{$\vect{F}_4$};

  % Resultant
  \draw[->, thick, codexgreen, line width=2.5pt, dashed]
    (O) -- (1.0,4.5)
    node[midway, left]{$\vect{R}$};

  \fill (O) circle (2.5pt) node[left]{Start};
  \fill (1.0,4.5) circle (2.5pt) node[above]{End};
  
  \node at (4.5,-0.9) {\small Resultant = tail of $\vect{F}_1$
    to head of $\vect{F}_4$};
\end{tikzpicture}
\end{center}
\end{diagbox}

\begin{diagbox}[Closed Polygon: Equilibrium]
\begin{center}
\begin{tikzpicture}[scale=0.85]
  \coordinate (O) at (0,0);
  \coordinate (A) at (3,0);
  \coordinate (B) at (4,2);
  \coordinate (C) at (1.5,3);

  \draw[->, thick, codexblue, line width=2pt]
    (O) -- (A) node[midway, below]{$\vect{F}_1$};
  \draw[->, thick, codexred, line width=2pt]
    (A) -- (B) node[midway, right]{$\vect{F}_2$};
  \draw[->, thick, codexgold, line width=2pt]
    (B) -- (C) node[midway, right]{$\vect{F}_3$};
  \draw[->, thick, violet, line width=2pt]
    (C) -- (O) node[midway, left]{$\vect{F}_4$};

  \node at (2,-0.9)
    {\small Polygon closes $\Rightarrow$
    $\vect{F}_1+\vect{F}_2+\vect{F}_3+\vect{F}_4 = \vect{0}$
    (Equilibrium)};
\end{tikzpicture}
\end{center}
\end{diagbox}

\noindent The polygon rule is the geometric basis
for the analytical component method developed in
Section~\ref{sec:component_method}. Instead of drawing
to scale, we replace geometry with arithmetic --- but
the logic is identical.

% ============================================================
\section{Calculating the Resultant Analytically}
% ============================================================

The graphical methods work, but for precise answers we use
trigonometry and the cosine rule.

\begin{lawbox}[Resultant of Two Vectors at Angle $\theta$]
If two vectors $\vect{A}$ and $\vect{B}$ are at an angle
$\theta$ to each other, the magnitude of their resultant is:
\[
  R = \sqrt{A^2 + B^2 + 2AB\cos\theta}
\]
The direction of the resultant (angle $\alpha$ from
$\vect{A}$) is:
\[
  \tan\alpha = \frac{B\sin\theta}{A + B\cos\theta}
\]
\end{lawbox}

\textbf{Special cases:}

\begin{itemize}
  \item $\theta = 0°$ (parallel, same direction):
  $R = A + B$ (maximum resultant)
  \item $\theta = 180°$ (antiparallel, opposite directions):
  $R = |A - B|$ (minimum resultant)
  \item $\theta = 90°$ (perpendicular):
  $R = \sqrt{A^2 + B^2}$, $\tan\alpha = B/A$
\end{itemize}

\begin{examplebox}[2.3]
\stepcounter{workedexample}
Two forces act on a bolt. Force $\vect{F}_1 = 80\ \text{N}$
acts horizontally to the right and force
$\vect{F}_2 = 60\ \text{N}$ acts vertically upward.
Find the magnitude and direction of the resultant force.
\end{examplebox}

\begin{solutionbox}
Since the two forces are perpendicular ($\theta = 90°$):
\[
  R = \sqrt{F_1^2 + F_2^2}
  = \sqrt{80^2 + 60^2}
  = \sqrt{6400 + 3600}
  = \sqrt{10000}
  = 100\ \text{N}
\]

Direction (angle above horizontal):
\[
  \tan\alpha = \frac{F_2}{F_1} = \frac{60}{80} = 0.75
\]
\[
  \alpha = \tan^{-1}(0.75) = 36.9°
\]

\textbf{Result:} The resultant force is
$\mathbf{100\ \text{N}}$ at $\mathbf{36.9°}$ above the
horizontal.

\begin{center}
\begin{tikzpicture}[scale=0.6]
  \draw[->, thick, codexblue, line width=2pt]
    (0,0) -- (4,0) node[right]{$F_1 = 80$ N};
  \draw[->, thick, codexred, line width=2pt]
    (0,0) -- (0,3) node[above]{$F_2 = 60$ N};
  \draw[->, thick, codexgreen, line width=2.5pt, dashed]
    (0,0) -- (4,3) node[right]{$R = 100$ N};
  \draw[gray, dashed] (4,0) -- (4,3);
  \draw[codexgold] (0.8,0) arc(0:36.9:0.8);
  \node[codexgold] at (1.2,0.3) {$36.9°$};
\end{tikzpicture}
\end{center}
\end{solutionbox}

\begin{examplebox}[2.4]
\stepcounter{workedexample}
Two tugboats pull a ship. Tugboat A pulls with a force of
$15\,000\ \text{N}$ and tugboat B pulls with a force of
$20\,000\ \text{N}$. The angle between the two tow ropes
is $60°$. Find the magnitude of the resultant force on
the ship.
\end{examplebox}

\begin{solutionbox}
Using the cosine rule with $A = 15\,000\ \text{N}$,
$B = 20\,000\ \text{N}$, $\theta = 60°$:
\begin{align*}
  R &= \sqrt{A^2 + B^2 + 2AB\cos\theta}\\[4pt]
  &= \sqrt{15000^2 + 20000^2
    + 2(15000)(20000)\cos 60°}\\[4pt]
  &= \sqrt{225{,}000{,}000 + 400{,}000{,}000
    + 600{,}000{,}000 \times 0.5}\\[4pt]
  &= \sqrt{225{,}000{,}000 + 400{,}000{,}000
    + 300{,}000{,}000}\\[4pt]
  &= \sqrt{925{,}000{,}000}\\[4pt]
  &= 30{,}414\ \text{N} \approx 30{,}400\ \text{N}
\end{align*}

\textbf{Result:} The resultant force on the ship is
approximately $\mathbf{30\,400\ \text{N}}$.

\begin{realworldbox}[Tugboats at Apapa Port]
At Apapa Port in Lagos, large vessels are routinely guided
into dock by multiple tugboats. The captains must account
for the vector sum of all tug forces, plus water current
and wind force, to bring a ship safely to berth.
Understanding resultant forces is not academic --- it
is how ships avoid hitting the dock.
\end{realworldbox}
\end{solutionbox}

% ============================================================
\section{Resolution of Vectors into Components}
% ============================================================

Just as we can \textit{add} vectors to find a resultant,
we can \textit{split} any vector into two perpendicular
components. This is called \textbf{resolution}.

Resolution is arguably more useful than addition, because
it allows us to handle any number of vectors by reducing
the problem to straightforward arithmetic.

\begin{lawbox}[Resolution of a Vector]
Any vector $\vect{F}$ at angle $\theta$ above the
horizontal can be resolved into:
\[
  F_x = F\cos\theta \quad \text{(horizontal component)}
\]
\[
  F_y = F\sin\theta \quad \text{(vertical component)}
\]
These components are perpendicular to each other and
together are completely equivalent to the original vector.
\end{lawbox}

\begin{diagbox}[Resolution of a Vector]
\begin{center}
\begin{tikzpicture}[scale=1.0]
  % Axes
  \draw[->, gray] (-0.3,0) -- (4.5,0) node[right]{$x$};
  \draw[->, gray] (0,-0.3) -- (0,3.5) node[above]{$y$};

  % Original vector
  \draw[->, thick, codexblue, line width=2.5pt]
    (0,0) -- (4,3)
    node[above right]{$\vect{F}$};

  % Components
  \draw[->, thick, codexred, line width=2pt]
    (0,0) -- (4,0)
    node[below]{$F_x = F\cos\theta$};
  \draw[->, thick, codexgreen, line width=2pt]
    (0,0) -- (0,3)
    node[left]{$F_y = F\sin\theta$};

  % Dashed completion
  \draw[dashed, gray] (4,0) -- (4,3);
  \draw[dashed, gray] (0,3) -- (4,3);

  % Angle
  \draw[codexgold] (0.8,0) arc(0:36.87:0.8);
  \node[codexgold] at (1.1,0.25) {$\theta$};

  % Right angle mark
  \draw[gray] (3.8,0) -- (3.8,0.2) -- (4,0.2);
\end{tikzpicture}
\end{center}
\end{diagbox}

\begin{examplebox}[2.5]
\stepcounter{workedexample}
A force of $250\ \text{N}$ acts at $35°$ above the
horizontal. Find the horizontal and vertical components.
\end{examplebox}

\begin{solutionbox}
\begin{align*}
  F_x &= F\cos\theta = 250\cos 35°
       = 250 \times 0.8192
       = 204.8\ \text{N}\\[6pt]
  F_y &= F\sin\theta = 250\sin 35°
       = 250 \times 0.5736
       = 143.4\ \text{N}
\end{align*}

\textbf{Check:} $\sqrt{F_x^2 + F_y^2}
= \sqrt{204.8^2 + 143.4^2}
= \sqrt{41\,943 + 20\,564}
= \sqrt{62\,507} \approx 250\ \text{N}$ ✓
\end{solutionbox}

% ============================================================
\section{Adding Multiple Vectors by Components}
\label{sec:component_method}
% ============================================================

When more than two vectors must be added, the most
powerful method is to resolve each vector into $x$ and $y$
components, sum the components separately, then combine
the sums.

\begin{lawbox}[Component Method for Multiple Vectors]
To find the resultant of vectors
$\vect{F}_1, \vect{F}_2, \ldots, \vect{F}_n$:
\begin{enumerate}[noitemsep]
  \item Resolve each vector into $x$ and $y$ components:
  $F_{ix} = F_i\cos\theta_i$,
  $F_{iy} = F_i\sin\theta_i$
  \item Sum all $x$-components: $R_x = \sum F_{ix}$
  \item Sum all $y$-components: $R_y = \sum F_{iy}$
  \item Find the resultant magnitude:
  $R = \sqrt{R_x^2 + R_y^2}$
  \item Find the resultant direction:
  $\alpha = \tan^{-1}\!\left(\dfrac{|R_y|}{|R_x|}\right)$
  (in the correct quadrant)
\end{enumerate}
\end{lawbox}

\begin{examplebox}[2.6]
\stepcounter{workedexample}
Three forces act on a particle simultaneously:
\begin{itemize}[noitemsep]
  \item $\vect{F}_1 = 40\ \text{N}$ at $0°$ (east)
  \item $\vect{F}_2 = 30\ \text{N}$ at $90°$ (north)
  \item $\vect{F}_3 = 50\ \text{N}$ at $210°$
  (southwest)
\end{itemize}
Find the magnitude and direction of the resultant.
\end{examplebox}

\begin{solutionbox}
\textbf{Step 1:} Resolve each force into components
(taking east as $+x$, north as $+y$):

\begin{center}
\begin{tabular}{lccc}
\toprule
\textbf{Force} & \textbf{Magnitude (N)} &
$\mathbf{F_x}$ \textbf{(N)} &
$\mathbf{F_y}$ \textbf{(N)}\\
\midrule
$\vect{F}_1$ & 40 &
  $40\cos 0° = +40.0$ &
  $40\sin 0° = 0.0$\\
$\vect{F}_2$ & 30 &
  $30\cos 90° = 0.0$ &
  $30\sin 90° = +30.0$\\
$\vect{F}_3$ & 50 &
  $50\cos 210° = -43.3$ &
  $50\sin 210° = -25.0$\\
\midrule
\textbf{Sum} & & $R_x = -3.3$ & $R_y = +5.0$\\
\bottomrule
\end{tabular}
\end{center}

\textbf{Step 2:} Find resultant magnitude:
\[
  R = \sqrt{(-3.3)^2 + (5.0)^2}
  = \sqrt{10.89 + 25.0}
  = \sqrt{35.89}
  = 5.99 \approx 6.0\ \text{N}
\]

\textbf{Step 3:} Find direction. Since $R_x < 0$ and
$R_y > 0$, the resultant is in the second quadrant
(northwest):
\[
  \alpha = \tan^{-1}\!\left(\frac{5.0}{3.3}\right)
  = \tan^{-1}(1.515) = 56.6°
\]

Angle from east (measuring anticlockwise):
$180° - 56.6° = 123.4°$

\textbf{Result:} $R \approx 6.0\ \text{N}$ at $123.4°$
from east (i.e.\ $56.6°$ west of north).
\end{solutionbox}

% ============================================================
\section{Vector Subtraction and Relative Velocity}
% ============================================================

\begin{definitionbox}[Vector Subtraction]
Subtracting vector $\vect{B}$ from $\vect{A}$ is equivalent
to adding the negative of $\vect{B}$:
\[
  \vect{A} - \vect{B} = \vect{A} + (-\vect{B})
\]
The negative of a vector has the same magnitude but
opposite direction.
\end{definitionbox}

\begin{definitionbox}[Relative Velocity]
The \textbf{velocity of A relative to B} is the velocity
of A as observed by B:
\[
  \vect{v}_{A\,\text{relative to}\,B}
  = \vect{v}_A - \vect{v}_B
\]
\end{definitionbox}

\begin{examplebox}[2.7]
\stepcounter{workedexample}
Car A travels due north at $60\ \text{km\,h}^{-1}$ and
car B travels due east at $80\ \text{km\,h}^{-1}$.
Find the velocity of car A as seen by a passenger
in car B.
\end{examplebox}

\begin{solutionbox}
\[
  \vect{v}_{A\,\text{rel}\,B}
  = \vect{v}_A - \vect{v}_B
\]

Taking north as $+y$ and east as $+x$:
\[
  \vect{v}_A = (0,\, 60)\ \text{km\,h}^{-1}
\]
\[
  \vect{v}_B = (80,\, 0)\ \text{km\,h}^{-1}
\]
\[
  \vect{v}_{A\,\text{rel}\,B}
  = (0 - 80,\; 60 - 0)
  = (-80,\, 60)\ \text{km\,h}^{-1}
\]

Magnitude:
\[
  |\vect{v}_{A\,\text{rel}\,B}|
  = \sqrt{80^2 + 60^2}
  = \sqrt{6400 + 3600} = 100\ \text{km\,h}^{-1}
\]

Direction: $\tan\alpha = 60/80 = 0.75$,
so $\alpha = 36.9°$ north of west.

\textbf{Result:} To the passenger in car B, car A appears
to move at $100\ \text{km\,h}^{-1}$ in a direction
$36.9°$ north of west.
\end{solutionbox}

\begin{realworldbox}[Danfo on Lagos Expressway]
When a danfo (minibus) overtakes you on the Lagos-Ibadan
Expressway, the relative velocity of the danfo with
respect to your car determines the rate at which it
appears to move past you. If you are doing 100 km\,h$^{-1}$
and the danfo is doing 120 km\,h$^{-1}$ in the same
direction, it appears to move past you at only
20 km\,h$^{-1}$ --- slow enough to read its side
advertisements. If it were coming toward you in the
opposite lane at 100 km\,h$^{-1}$, the relative velocity
would be 200 km\,h$^{-1}$ --- which is why head-on
collisions are so catastrophic.
\end{realworldbox}

% ============================================================
\section{Position Vectors}
\label{sec:position_vectors}
% ============================================================

\textbf{Intuition:} A position vector answers a very
specific question: \textit{where is this point relative
to the origin?} It is an arrow drawn from the fixed
reference point $O$ (the origin) to the point in
question.

\begin{definitionbox}[Position Vector]
The \textbf{position vector} $\vect{r}_P$ of a point
$P(x,\,y)$ is:
\[
  \vect{r}_P = x\hat{i} + y\hat{j}
  \quad \text{(2D)}
\]
\[
  \vect{r}_P = x\hat{i} + y\hat{j} + z\hat{k}
  \quad \text{(3D)}
\]
Its magnitude $|\vect{r}_P|$ is the distance from the
origin to $P$.
\end{definitionbox}

\begin{diagbox}[Position Vectors in 2D]
\begin{center}
\begin{tikzpicture}[scale=0.8]
  % Axes
  \draw[->, gray] (-0.5,0) -- (5.5,0) node[right]{$x$};
  \draw[->, gray] (0,-0.5) -- (0,4.0) node[above]{$y$};

  % Point A
  \coordinate (A) at (1.5,3.0);
  \fill[codexblue] (A) circle (3pt) node[above right]{$A(1.5,\,3.0)$};
  \draw[->, thick, codexblue, line width=2pt]
    (0,0) -- (A) node[midway, left]{$\vect{r}_A$};

  % Point B
  \coordinate (B) at (4.5,1.5);
  \fill[codexred] (B) circle (3pt) node[above right]{$B(4.5,\,1.5)$};
  \draw[->, thick, codexred, line width=2pt]
    (0,0) -- (B) node[midway, below]{$\vect{r}_B$};

  % Vector from A to B
  \draw[->, thick, codexgreen, line width=2pt, dashed]
    (A) -- (B) node[pos=0.55, below=11pt, fill=white, inner sep=1pt]{$\vect{AB}$};

  \fill (0,0) circle (3pt) node[below left]{$O$};
\end{tikzpicture}
\end{center}
\end{diagbox}

% ----------------------------------------------------------
\subsection*{Vector Between Two Points}

To find the vector pointing from point $A$ to point $B$,
subtract the position vector of $A$ from that of $B$:

\begin{lawbox}[Vector Between Two Points]
\[
  \vect{AB} = \vect{r}_B - \vect{r}_A
\]
This gives the displacement vector from $A$ to $B$.
Its magnitude $|\vect{AB}|$ is the distance between $A$ and $B$.
\end{lawbox}

\begin{examplebox}[2.8]
\stepcounter{workedexample}
Point $A$ has coordinates $(2,\,5)$ and point $B$ has
coordinates $(7,\,2)$ (in metres).\\
(a) Write down the position vectors $\vect{r}_A$ and
$\vect{r}_B$.\\
(b) Find the vector $\vect{AB}$ from $A$ to $B$.\\
(c) Find the distance $|AB|$.
\end{examplebox}

\begin{solutionbox}
(a) Position vectors:
\[
  \vect{r}_A = 2\hat{i} + 5\hat{j}\ \text{m},\qquad
  \vect{r}_B = 7\hat{i} + 2\hat{j}\ \text{m}
\]

(b) Vector from $A$ to $B$:
\[
  \vect{AB} = \vect{r}_B - \vect{r}_A
  = (7-2)\hat{i} + (2-5)\hat{j}
  = 5\hat{i} - 3\hat{j}\ \text{m}
\]

(c) Distance:
\[
  |AB| = |\vect{AB}| = \sqrt{5^2 + (-3)^2}
  = \sqrt{25 + 9} = \sqrt{34} \approx 5.83\ \text{m}
\]

\textbf{Note:} $\vect{BA} = -\vect{AB} = -5\hat{i} + 3\hat{j}$
--- same distance, opposite direction.
\end{solutionbox}

% ============================================================
\section{Unit Vectors}
\label{sec:unit_vectors}
% ============================================================

\begin{definitionbox}[Unit Vector]
A \textbf{unit vector} is a vector with magnitude exactly
equal to 1. Its sole purpose is to specify a direction.
The standard unit vectors in three dimensions are:
\[
  \hat{i} \text{ (east/x-direction)},\quad
  \hat{j} \text{ (north/y-direction)},\quad
  \hat{k} \text{ (upward/z-direction)}
\]
Any vector can be written as:
$\vect{F} = F_x\hat{i} + F_y\hat{j} + F_z\hat{k}$
\end{definitionbox}

\begin{examplebox}[2.9]
\stepcounter{workedexample}
A displacement is $\vect{s} = 3\hat{i} + 4\hat{j}$
metres. Find its magnitude and direction.
\end{examplebox}

\begin{solutionbox}
Magnitude:
\[
  s = |\vect{s}| = \sqrt{3^2 + 4^2}
  = \sqrt{9 + 16} = \sqrt{25} = 5\ \text{m}
\]
Direction:
\[
  \theta = \tan^{-1}\!\left(\frac{4}{3}\right)
  = \tan^{-1}(1.333) = 53.1°
  \text{ above the } x\text{-axis (north of east)}
\]
\end{solutionbox}

% ----------------------------------------------------------
\subsection*{Constructing a Unit Vector from Any Vector}

\textbf{Intuition:} To extract the direction from a
vector while discarding its size, divide every component
by the magnitude. This scales the arrow down to length 1
without rotating it.

\begin{lawbox}[Unit Vector from a Given Vector]
The unit vector in the direction of $\vect{A}$ is:
\[
  \hat{A} = \frac{\vect{A}}{|\vect{A}|}
\]
\textbf{Verification:} $|\hat{A}|
= \dfrac{|\vect{A}|}{|\vect{A}|} = 1$ always. $\checkmark$
\end{lawbox}

\begin{examplebox}[2.10]
\stepcounter{workedexample}
Find the unit vector in the direction of
$\vect{F} = 6\hat{i} + 8\hat{j}\ \text{N}$.
\end{examplebox}

\begin{solutionbox}
\textbf{Step 1:} Find the magnitude:
\[
  |\vect{F}| = \sqrt{6^2 + 8^2}
  = \sqrt{36 + 64} = \sqrt{100} = 10\ \text{N}
\]

\textbf{Step 2:} Divide by the magnitude:
\[
  \hat{F} = \frac{\vect{F}}{|\vect{F}|}
  = \frac{6\hat{i} + 8\hat{j}}{10}
  = 0.6\hat{i} + 0.8\hat{j}
\]

\textbf{Verification:}
$|\hat{F}| = \sqrt{0.6^2 + 0.8^2}
= \sqrt{0.36 + 0.64} = \sqrt{1.00} = 1$ ✓
\end{solutionbox}

\begin{examplebox}[2.11]
\stepcounter{workedexample}
Find the unit vector in the direction of
$\vect{v} = 2\hat{i} - 3\hat{j} + 6\hat{k}\ \text{m\,s}^{-1}$.
\end{examplebox}

\begin{solutionbox}
\textbf{Step 1:} Magnitude (from Example~2.2):
\[
  |\vect{v}| = \sqrt{4 + 9 + 36} = 7\ \text{m\,s}^{-1}
\]

\textbf{Step 2:} Unit vector:
\[
  \hat{v} = \frac{2\hat{i} - 3\hat{j} + 6\hat{k}}{7}
  = \frac{2}{7}\hat{i} - \frac{3}{7}\hat{j}
  + \frac{6}{7}\hat{k}
\]

\textbf{Verification:}
$|\hat{v}| = \sqrt{(2/7)^2+(-3/7)^2+(6/7)^2}
= \sqrt{4+9+36}/7 = 7/7 = 1$ ✓
\end{solutionbox}

% ============================================================
\section{Equilibrium of Forces}
% ============================================================

\begin{definitionbox}[Equilibrium]
A particle is in \textbf{equilibrium} when the resultant
of all forces acting on it is zero:
\[
  \sum \vect{F} = \vect{0}
\]
This means $\sum F_x = 0$ AND $\sum F_y = 0$
simultaneously.
\end{definitionbox}

\begin{lawbox}[Lami's Theorem]
If three concurrent coplanar forces $F_1$, $F_2$, $F_3$
are in equilibrium, then each force is proportional to
the sine of the angle between the other two:
\[
  \frac{F_1}{\sin\alpha_1}
  = \frac{F_2}{\sin\alpha_2}
  = \frac{F_3}{\sin\alpha_3}
\]
where $\alpha_1$, $\alpha_2$, $\alpha_3$ are the angles
\textit{opposite} to $F_1$, $F_2$, $F_3$ respectively.
\end{lawbox}

\begin{examplebox}[2.12]
\stepcounter{workedexample}
A traffic light of weight $W = 400\ \text{N}$ is
suspended from two cables. Cable 1 makes an angle of
$30°$ with the vertical and cable 2 makes an angle of
$45°$ with the vertical. Find the tension in each cable.
\end{examplebox}

\begin{solutionbox}
\begin{center}
\begin{tikzpicture}[scale=0.8]
  % Ceiling
  \fill[gray!30] (-3,2) rectangle (3,2.3);
  \draw[thick] (-3,2) -- (3,2);
  % Cables
  % 30 degrees from the vertical: |x| = 2 tan(30 degrees) = 1.155.
  \draw[thick, codexblue] (0,0) -- (-1.155, 2)
    node[midway, left]{$T_1$};
  \draw[thick, codexred] (0,0) -- (2,2)
    node[midway, right]{$T_2$};
  % Weight
  \draw[->, thick, codexgreen, line width=2pt]
    (0,0) -- (0,-1.5) node[below]{$W = 400$ N};
  % Angles
  \draw[dashed, gray] (0,0) -- (0,2.5);
  \draw[codexblue] (0,0.5) arc(90:120:0.5);
  \node[left, codexblue, font=\small] at (-0.65,0.55)
    {$30°$};
  \draw[codexred] (0,0.5) arc(90:45:0.5);
  \node[right, codexred, font=\small] at (0.0,0.8)
    {$45°$};
  \fill (0,0) circle (3pt);
\end{tikzpicture}
\end{center}

For equilibrium, resolve horizontally and vertically:
\begin{align*}
  \sum F_x = 0&:\quad
  T_2\sin 45° - T_1\sin 30° = 0\\
  \sum F_y = 0&:\quad
  T_1\cos 30° + T_2\cos 45° - 400 = 0
\end{align*}

From the first equation:
\[
  T_2 = T_1 \frac{\sin 30°}{\sin 45°}
  = T_1 \times \frac{0.500}{0.707} = 0.707 T_1
\]

Substituting into the second equation:
\[
  T_1(0.866) + 0.707 T_1(0.707) = 400
\]
\[
  0.866 T_1 + 0.500 T_1 = 400
\]
\[
  1.366 T_1 = 400
\]
\[
  \boxed{T_1 = 293\ \text{N}}
\]
\[
  \boxed{T_2 = 0.707 \times 293 = 207\ \text{N}}
\]
\end{solutionbox}

% ============================================================
\section{The Scalar (Dot) Product}
\label{sec:dot_product}
% ============================================================

\textbf{Intuition:} Imagine two forces acting on a box
at an angle to each other. You want to know how much
of one force acts \textit{in the direction} of the
other --- i.e.\ how much of $\vect{A}$ is projected
onto $\vect{B}$. That projection, multiplied by the
magnitude of $\vect{B}$, is the dot product.

More concretely: if you push a box at an angle to its
direction of motion, only the component of your push
\textit{along} the motion does useful work. Work is
the classic physical realisation of the dot product.

\begin{definitionbox}[Scalar (Dot) Product]
The \textbf{scalar product} (also called the
\textbf{dot product}) of two vectors $\vect{A}$ and
$\vect{B}$ is defined as:
\[
  \vect{A} \cdot \vect{B}
  = AB\cos\theta
\]
where $A = |\vect{A}|$, $B = |\vect{B}|$, and $\theta$
is the angle between the two vectors
($0° \leq \theta \leq 180°$).

The result is a \textbf{scalar} --- a single number
with no direction.
\end{definitionbox}

\begin{diagbox}[Geometric Meaning of the Dot Product]
\begin{center}
\begin{tikzpicture}[scale=0.9]
  % Vectors
  \draw[->, thick, codexblue, line width=2pt]
    (0,0) -- (4,0) node[right]{$\vect{B}$};
  \draw[->, thick, codexred, line width=2pt]
    (0,0) -- (3.5,2.5) node[above right]{$\vect{A}$};

  % Angle
  \draw[codexgold] (0.9,0) arc(0:35.5:0.9);
  \node[codexgold] at (1.3,0.3) {$\theta$};

  % Projection of A onto B
  \draw[dashed, gray] (3.5,2.5) -- (3.5,0);
  % Offset the projection arrow below B's axis so both vectors are visible.
  \draw[densely dotted, gray] (0,0) -- (0,-0.35);
  \draw[densely dotted, gray] (3.5,0) -- (3.5,-0.35);
  \draw[->, thick, codexgreen, line width=2pt]
    (0,-0.35) -- (3.5,-0.35)
    node[below]{$A\cos\theta$};

  \node at (2,-1.2)
    {\small $\vect{A}\cdot\vect{B}
    = B \times (A\cos\theta)
    = $ (magnitude of $\vect{B}$)
    $\times$ (projection of $\vect{A}$ onto $\vect{B}$)};
\end{tikzpicture}
\end{center}
\end{diagbox}

% ----------------------------------------------------------
\subsection*{Component Form}

When vectors are expressed in component form, the dot
product is computed without needing the angle:

\begin{lawbox}[Dot Product in Component Form]
\textbf{2D:}
$\vect{A}\cdot\vect{B}
= A_x B_x + A_y B_y$

\vspace{0.2cm}

\textbf{3D:}
$\vect{A}\cdot\vect{B}
= A_x B_x + A_y B_y + A_z B_z$
\end{lawbox}

\textit{Why does this work?} Expand using the
distributive law and the unit vector dot products
($\hat{i}\cdot\hat{i} = 1$,
$\hat{i}\cdot\hat{j} = 0$, etc.) and you recover
the component formula directly.

% ----------------------------------------------------------
\subsection*{Properties and Special Results}

\begin{keybox}[Unit Vector Dot Products]
\centering
\begin{tabular}{ccccc}
$\hat{i}\cdot\hat{i}=1$ &
$\hat{j}\cdot\hat{j}=1$ &
$\hat{k}\cdot\hat{k}=1$ \\[4pt]
$\hat{i}\cdot\hat{j}=0$ &
$\hat{j}\cdot\hat{k}=0$ &
$\hat{k}\cdot\hat{i}=0$ \\
\end{tabular}

\vspace{0.2cm}
\textit{Same basis vectors} $\to$ dot product = 1.
\textit{Different basis vectors} $\to$ dot product = 0.
\end{keybox}

\textbf{Properties:}
\begin{align*}
  \vect{A}\cdot\vect{B}
  &= \vect{B}\cdot\vect{A}
  \quad \text{(commutative)}\\
  \vect{A}\cdot(\vect{B}+\vect{C})
  &= \vect{A}\cdot\vect{B}
  + \vect{A}\cdot\vect{C}
  \quad \text{(distributive)}\\
  \vect{A}\cdot\vect{A}
  &= A^2
  \quad \text{(magnitude squared)}
\end{align*}

\textbf{Special cases:}
\begin{itemize}[noitemsep]
  \item \textbf{Perpendicular vectors:}
  $\vect{A}\cdot\vect{B} = 0$
  (since $\cos 90° = 0$) ---
  use this to \textit{test} for perpendicularity.
  \item \textbf{Parallel vectors:}
  $\vect{A}\cdot\vect{B} = AB$
  (since $\cos 0° = 1$).
  \item \textbf{Anti-parallel vectors:}
  $\vect{A}\cdot\vect{B} = -AB$
  (since $\cos 180° = -1$).
\end{itemize}

% ----------------------------------------------------------
\subsection*{Finding the Angle Between Two Vectors}

Rearranging the definition $\vect{A}\cdot\vect{B}
= AB\cos\theta$ gives a powerful formula:

\begin{lawbox}[Angle Between Two Vectors]
\[
  \theta
  = \cos^{-1}\!\left(
    \frac{\vect{A}\cdot\vect{B}}{|\vect{A}||\vect{B}|}
  \right)
\]
This formula works in both 2D and 3D.
\end{lawbox}

% ----------------------------------------------------------
\subsection*{Physical Application: Work}

\begin{keybox}[Work as a Dot Product]
When a constant force $\vect{F}$ displaces an object
by $\vect{s}$, the work done is:
\[
  W = \vect{F}\cdot\vect{s} = Fs\cos\theta
\]
Only the component of force \textit{along} the
displacement does work. If $\theta = 90°$, no
work is done --- regardless of how large the force is.
\end{keybox}

\begin{examplebox}[2.13]
\stepcounter{workedexample}
Two forces are
$\vect{F}_1 = 3\hat{i} + 4\hat{j}\ \text{N}$ and
$\vect{F}_2 = 5\hat{i} - 3\hat{j}\ \text{N}$.\\
(a) Find $\vect{F}_1 \cdot \vect{F}_2$.\\
(b) Find the angle between the two forces.\\
(c) Are the forces perpendicular?
\end{examplebox}

\begin{solutionbox}
(a) Component form:
\[
  \vect{F}_1 \cdot \vect{F}_2
  = (3)(5) + (4)(-3)
  = 15 - 12 = 3\ \text{N}^2
\]

(b) Magnitudes:
\[
  |\vect{F}_1| = \sqrt{9+16} = 5\ \text{N},\qquad
  |\vect{F}_2| = \sqrt{25+9} = \sqrt{34}\ \text{N}
\]
Angle:
\[
  \theta = \cos^{-1}\!\left(\frac{3}{5\sqrt{34}}\right)
  = \cos^{-1}\!\left(\frac{3}{29.15}\right)
  = \cos^{-1}(0.1029)
  = 84.1°
\]

(c) Since $\vect{F}_1 \cdot \vect{F}_2 = 3 \neq 0$,
the forces are \textbf{not} perpendicular.
\end{solutionbox}

\begin{examplebox}[2.14]
\stepcounter{workedexample}
A force $\vect{F} = 4\hat{i} + 3\hat{j}\ \text{N}$
acts on an object that undergoes a displacement
$\vect{s} = 5\hat{i} + 2\hat{j}\ \text{m}$.
Find the work done by the force.
\end{examplebox}

\begin{solutionbox}
\[
  W = \vect{F}\cdot\vect{s}
  = (4)(5) + (3)(2)
  = 20 + 6
  = 26\ \text{J}
\]

\textbf{Note:} Using $W = Fs\cos\theta$ would give
the same answer but requires more steps. The component
formula is faster when vectors are already in
component form.
\end{solutionbox}

% ============================================================
\section{The Vector (Cross) Product}
\label{sec:cross_product}
% ============================================================

\textbf{Intuition:} The cross product asks a different
question to the dot product. Instead of ``how aligned
are these two vectors?'', it asks ``how
\textit{perpendicular} are they?'' and produces a
brand-new vector that points \textit{out of} the plane
containing the original two.

Think of tightening a bolt with a spanner. The force
on the spanner and the arm of the spanner are both
vectors. The \textit{torque} (turning effect) is
perpendicular to both --- it points along the axis
of rotation. That perpendicular vector is the cross
product.

\begin{definitionbox}[Vector (Cross) Product]
The \textbf{vector product} (also called the
\textbf{cross product}) of $\vect{A}$ and $\vect{B}$
is a vector $\vect{C}$ defined by:
\[
  \vect{C} = \vect{A}\times\vect{B}
\]
\[
  |\vect{C}| = AB\sin\theta
\]
where $\theta$ is the angle between $\vect{A}$ and
$\vect{B}$ ($0° \leq \theta \leq 180°$).

The \textbf{direction} of $\vect{C}$ is perpendicular
to the plane containing $\vect{A}$ and $\vect{B}$,
determined by the right-hand rule.

The result is a \textbf{vector}.
\end{definitionbox}

\begin{keybox}[Geometric Meaning]
$|\vect{A}\times\vect{B}| = AB\sin\theta$
is the \textbf{area of the parallelogram} formed by
$\vect{A}$ and $\vect{B}$ as adjacent sides.
\end{keybox}

% ----------------------------------------------------------
\subsection*{The Right-Hand Rule}

The right-hand rule gives the direction of
$\vect{A}\times\vect{B}$:

\begin{lawbox}[Right-Hand Rule]
\begin{enumerate}[noitemsep]
  \item Point the fingers of your right hand in the
  direction of $\vect{A}$.
  \item Curl them towards $\vect{B}$ (through the
  smaller angle $\theta$).
  \item Your extended right thumb points in the
  direction of $\vect{A}\times\vect{B}$.
\end{enumerate}
\end{lawbox}

\begin{diagbox}[Right-Hand Rule]
\begin{center}
\begin{tikzpicture}[scale=0.9]
  % Draw A and B in xy-plane
  \draw[->, thick, codexblue, line width=2pt]
    (0,0) -- (3,0) node[right]{$\vect{A}$};
  \draw[->, thick, codexred, line width=2pt]
    (0,0) -- (1.5,2.5) node[above right]{$\vect{B}$};

  % Cross product direction: circled dot denotes out of the page.
  \node[draw, circle, thick, codexgreen,
    inner sep=3pt] at (-1.5,1.5) {$\cdot$};
  \node[right, codexgreen, font=\small] at (-1.25,1.5)
    {$\vect{A}\times\vect{B}$ (out of page)};

  \draw[codexgold] (0.7,0) arc(0:59:0.7);
  \node[codexgold] at (1.0,0.35) {$\theta$};

  \node at (1.5,-0.6)
    {\small Right-hand fingers: $\vect{A}\to\vect{B}$,
    thumb: $\vect{A}\times\vect{B}$};
\end{tikzpicture}
\end{center}
\end{diagbox}

% ----------------------------------------------------------
\subsection*{Anti-commutativity}

\begin{lawbox}[Anti-commutativity of Cross Product]
\[
  \vect{A}\times\vect{B}
  = -(\vect{B}\times\vect{A})
\]
Reversing the order flips the direction of the result.
The cross product is \textbf{not} commutative.
\end{lawbox}

This makes physical sense: if $\vect{A}\times\vect{B}$
points upward (out of the page), then
$\vect{B}\times\vect{A}$ points downward (into the page).

% ----------------------------------------------------------
\subsection*{Unit Vector Cross Products: The Cyclic Rule}

\begin{keybox}[Cyclic Rule for Unit Vectors]
\centering
\begin{tikzpicture}[scale=0.8]
  \node[draw, circle, codexblue, thick,
    inner sep=4pt] (i) at (0,0) {$\hat{k}$};
  \node[draw, circle, codexred, thick,
    inner sep=4pt] (j) at (2.5,0) {$\hat{j}$};
  \node[draw, circle, codexgreen, thick,
    inner sep=4pt] (k) at (1.25,2) {$\hat{i}$};

  \draw[<-, thick, codexblue] (i) -- (j);
  \draw[<-, thick, codexred]  (j) -- (k);
  \draw[<-, thick, codexgreen](k) -- (i);

  \node at (1.25,-1.0)
    {\small Clockwise cycle: positive product};
    
\end{tikzpicture}

\vspace{0.3cm}
\begin{tabular}{lll}
$\hat{i}\times\hat{j} = \hat{k}$ &
$\hat{j}\times\hat{k} = \hat{i}$ &
$\hat{k}\times\hat{i} = \hat{j}$ \\[4pt]
$\hat{j}\times\hat{i} = -\hat{k}$ &
$\hat{k}\times\hat{j} = -\hat{i}$ &
$\hat{i}\times\hat{k} = -\hat{j}$ \\[4pt]
\multicolumn{3}{l}{$\hat{i}\times\hat{i}
= \hat{j}\times\hat{j}
= \hat{k}\times\hat{k} = \vect{0}$
\quad (parallel $\Rightarrow$ $\sin 0° = 0$)}
\end{tabular}
\end{keybox}

% ----------------------------------------------------------
\subsection*{Component Form via the Determinant}

For vectors in component form, compute the cross product
using the $3\times 3$ determinant expansion:

\vspace{2.0cm}
\begin{lawbox}[Cross Product in Component Form]
\[
  \vect{A}\times\vect{B}
  = \begin{vmatrix}
    \hat{i} & \hat{j} & \hat{k}\\
    A_x & A_y & A_z\\
    B_x & B_y & B_z
  \end{vmatrix}
\]

Expanding along the top row:
\begin{align*}
  \vect{A}\times\vect{B}
  =\;
  &(A_y B_z - A_z B_y)\,\hat{i}\\
  -\;&(A_x B_z - A_z B_x)\,\hat{j}\\
  +\;&(A_x B_y - A_y B_x)\,\hat{k}
\end{align*}

\textit{Memory aid:} For each component, multiply the
two remaining components diagonally and subtract
(``cross-multiply and subtract'').
\end{lawbox}

% ----------------------------------------------------------
\subsection*{Special Cases}

\begin{itemize}[noitemsep]
  \item \textbf{Parallel vectors} ($\theta = 0°$ or
  $180°$): $\vect{A}\times\vect{B} = \vect{0}$
  (since $\sin 0° = \sin 180° = 0$) ---
  use this to \textit{test} for parallelism.
  \item \textbf{Perpendicular vectors} ($\theta = 90°$):
  $|\vect{A}\times\vect{B}| = AB$ (maximum cross product
  magnitude).
\end{itemize}

% ----------------------------------------------------------
\subsection*{Physical Applications}

\begin{keybox}[Torque as a Cross Product]
The torque (moment) $\vect{\tau}$ produced by a force
$\vect{F}$ applied at position $\vect{r}$ from the
pivot is:
\[
  \vect{\tau} = \vect{r}\times\vect{F}
\]
$|\vect{\tau}| = rF\sin\theta$, where $\theta$ is the
angle between $\vect{r}$ and $\vect{F}$.
The direction of $\vect{\tau}$ (given by the right-hand
rule) is along the axis of rotation.
\end{keybox}

\begin{keybox}[Magnetic Force as a Cross Product]
The force on a charge $q$ moving with velocity
$\vect{v}$ in a magnetic field $\vect{B}$ is:
\[
  \vect{F} = q\vect{v}\times\vect{B}
\]
This force is always perpendicular to both $\vect{v}$
and $\vect{B}$, which is why magnetic fields curve
moving charges into circular paths.
\end{keybox}

\begin{examplebox}[2.15]
\stepcounter{workedexample}
Find the cross product $\vect{A}\times\vect{B}$ where
$\vect{A} = 3\hat{i} + 2\hat{j} - \hat{k}$ and
$\vect{B} = \hat{i} - 3\hat{j} + 2\hat{k}$.
Also find the magnitude of the cross product.
\end{examplebox}

\begin{solutionbox}
Applying the determinant formula:
\begin{align*}
  \vect{A}\times\vect{B}
  &= \bigl[(2)(2) - (-1)(-3)\bigr]\hat{i}
   - \bigl[(3)(2) - (-1)(1)\bigr]\hat{j}
   + \bigl[(3)(-3) - (2)(1)\bigr]\hat{k}\\[6pt]
  &= [4 - 3]\hat{i}
   - [6 - (-1)]\hat{j}
   + [-9 - 2]\hat{k}\\[6pt]
  &= \hat{i} - 7\hat{j} - 11\hat{k}
\end{align*}

Magnitude:
\[
  |\vect{A}\times\vect{B}|
  = \sqrt{1^2 + (-7)^2 + (-11)^2}
  = \sqrt{1 + 49 + 121}
  = \sqrt{171}
  \approx 13.1
\]

\textbf{Verification (perpendicularity check):}
$(\vect{A}\times\vect{B})\cdot\vect{A}
= (1)(3)+(-7)(2)+(-11)(-1)
= 3 - 14 + 11 = 0$ ✓\\
$(\vect{A}\times\vect{B})\cdot\vect{B}
= (1)(1)+(-7)(-3)+(-11)(2)
= 1 + 21 - 22 = 0$ ✓

The result is perpendicular to both original vectors,
as expected.
\end{solutionbox}

\begin{examplebox}[2.16]
\stepcounter{workedexample}
A spanner applies a force $\vect{F} = 20\hat{j}\ \text{N}$
at a position $\vect{r} = 0.3\hat{i}\ \text{m}$ from
the pivot. Find the torque vector.
\end{examplebox}

\begin{solutionbox}
\[
  \vect{\tau} = \vect{r}\times\vect{F}
  = (0.3\hat{i})\times(20\hat{j})
  = (0.3)(20)(\hat{i}\times\hat{j})
  = 6\hat{k}\ \text{N\,m}
\]

The torque is $6\ \text{N\,m}$ in the $+\hat{k}$ direction
(i.e.\ upward, out of the page), corresponding to an
anticlockwise rotation. This matches the physical
expectation: pushing up on a horizontal spanner arm
produces an anticlockwise torque about the bolt.
\end{solutionbox}

% ============================================================
\section{Dot Product vs.\ Cross Product: A Decision Guide}
\label{sec:dot_vs_cross}
% ============================================================

Students sometimes confuse when to apply each product.
The key is to ask: \textit{what type of answer does the
physics require?}

\begin{keybox}[Dot vs.\ Cross Product Comparison]
\centering
\renewcommand{\arraystretch}{1.3}
\begin{tabular}{lll}
\toprule
\textbf{Feature} &
\textbf{Dot Product} $\vect{A}\cdot\vect{B}$ &
\textbf{Cross Product} $\vect{A}\times\vect{B}$\\
\midrule
Result type & Scalar & Vector\\
Formula & $AB\cos\theta$ & $AB\sin\theta$\\
Zero when & Perpendicular ($\theta=90°$)
  & Parallel ($\theta=0°$ or $180°$)\\
Maximum when & Parallel ($\theta=0°$)
  & Perpendicular ($\theta=90°$)\\
Commutative? & Yes: $\vect{A}\cdot\vect{B}
  = \vect{B}\cdot\vect{A}$
  & No: $\vect{A}\times\vect{B}
  = -\vect{B}\times\vect{A}$\\
Measures & Alignment / projection
  & Perpendicularity / rotation\\
Physics use & Work, projections, angles
  & Torque, magnetic force, area\\
\bottomrule
\end{tabular}
\end{keybox}

\textbf{Decision logic:}
\begin{itemize}[noitemsep]
  \item If the answer is a \textit{scalar} (a number:
  work done, power, angle between vectors) $\to$
  use the \textbf{dot product}.
  \item If the answer is a \textit{vector} perpendicular
  to both inputs (torque, magnetic force, angular
  momentum) $\to$ use the \textbf{cross product}.
  \item If you need to \textit{test for perpendicularity}
  $\to$ compute the dot product and check for zero.
  \item If you need to \textit{test for parallelism}
  $\to$ compute the cross product and check for
  the zero vector.
\end{itemize}

\vspace{3.5cm}

% ============================================================
\section{Chapter Summary}
% ============================================================

\begin{summarybox}

\textbf{Scalar:} Magnitude and unit only. Adds algebraically.

\textbf{Vector:} Magnitude, unit, and direction. Adds
by vector laws.

\vspace{0.2cm}

\textbf{Types of vectors:} Equal, null, negative,
parallel, anti-parallel, unit, position, free,
localised, collinear, coplanar.

\vspace{0.2cm}

\textbf{Magnitude:}
2D: $|\vect{A}| = \sqrt{A_x^2+A_y^2}$;\quad
3D: $|\vect{A}| = \sqrt{A_x^2+A_y^2+A_z^2}$

\textbf{Range of resultant:}
$|A-B| \leq R \leq A+B$

\vspace{0.2cm}

\textbf{Vector addition:}
\begin{itemize}[noitemsep]
  \item Triangle (head-to-tail) rule
  \item Parallelogram rule
  \item Polygon rule (extended head-to-tail; closed
  polygon $\Rightarrow$ equilibrium)
  \item Analytical (cosine rule):
  $R = \sqrt{A^2+B^2+2AB\cos\theta}$
\end{itemize}

\textbf{Resolution:}
$F_x = F\cos\theta$, $F_y = F\sin\theta$

\textbf{Component method:}
$R_x = \sum F_{ix}$, $R_y = \sum F_{iy}$,
$R = \sqrt{R_x^2 + R_y^2}$

\textbf{Vector subtraction:}
$\vect{A}-\vect{B} = \vect{A}+(-\vect{B})$

\textbf{Relative velocity:}
$\vect{v}_{A\,\text{rel}\,B} = \vect{v}_A - \vect{v}_B$

\textbf{Position vector:}
$\vect{r}_P = x\hat{i}+y\hat{j}$;\quad
$\vect{AB} = \vect{r}_B - \vect{r}_A$

\textbf{Unit vector from $\vect{A}$:}
$\hat{A} = \dfrac{\vect{A}}{|\vect{A}|}$,\quad
$|\hat{A}|=1$ always.

\textbf{Equilibrium:}
$\sum F_x = 0$ and $\sum F_y = 0$ simultaneously.

\textbf{Dot product:}
$\vect{A}\cdot\vect{B} = AB\cos\theta
= A_xB_x + A_yB_y + A_zB_z$
(scalar; zero when perpendicular)

\textbf{Angle between vectors:}
$\theta = \cos^{-1}\!\left(\dfrac{\vect{A}\cdot\vect{B}}{AB}\right)$

\textbf{Cross product:}
$|\vect{A}\times\vect{B}| = AB\sin\theta$;
direction by right-hand rule
(vector; zero when parallel)

\textbf{Cyclic rule:}
$\hat{i}\times\hat{j}=\hat{k}$,
$\hat{j}\times\hat{k}=\hat{i}$,
$\hat{k}\times\hat{i}=\hat{j}$

\textbf{Torque:} $\vect{\tau}=\vect{r}\times\vect{F}$;\quad
\textbf{Work:} $W=\vect{F}\cdot\vect{s}$

\end{summarybox}

% ============================================================
% EXERCISES
% ============================================================
\newpage

\begin{exercisebox}[2]

\subsection*{Section A: Multiple Choice}

\textbf{A1.} Which pair consists of two vector quantities?

\mcoption{A}{Speed and velocity}
\mcoption{B}{Displacement and velocity}
\mcoption{C}{Distance and mass}
\mcoption{D}{Power and force}
\ansref{Ch2-A1}

\vspace{0.4cm}

\textbf{A2.} A force of $8\ \text{N}$ and a force of
$6\ \text{N}$ act on a particle. Which of the following
cannot be the magnitude of the resultant?

\mcoption{A}{$2\ \text{N}$}
\mcoption{B}{$8\ \text{N}$}
\mcoption{C}{$14\ \text{N}$}
\mcoption{D}{$15\ \text{N}$}
\ansref{Ch2-A2}

\vspace{0.4cm}

\textbf{A3.} A vector $\vect{F}$ of magnitude $100\ \text{N}$
acts at $60°$ to the horizontal. The horizontal component
of $\vect{F}$ is:

\mcoption{A}{$50\ \text{N}$}
\mcoption{B}{$57.7\ \text{N}$}
\mcoption{C}{$86.6\ \text{N}$}
\mcoption{D}{$100\ \text{N}$}
\ansref{Ch2-A3}

\vspace{0.4cm}

\textbf{A4.} Three forces of $5\ \text{N}$, $12\ \text{N}$
and $13\ \text{N}$ act at a point. For these forces to be
in equilibrium, the angle between the $5\ \text{N}$ and
$12\ \text{N}$ forces must be:

\mcoption{A}{$30°$}
\mcoption{B}{$60°$}
\mcoption{C}{$90°$}
\mcoption{D}{$120°$}
\ansref{Ch2-A4}

\vspace{0.4cm}

\textbf{A5.} A vector has components
$F_x = -3\ \text{N}$ and $F_y = -4\ \text{N}$.
Its magnitude is:

\mcoption{A}{$1\ \text{N}$}
\mcoption{B}{$5\ \text{N}$}
\mcoption{C}{$7\ \text{N}$}
\mcoption{D}{$-5\ \text{N}$}
\ansref{Ch2-A5}

\vspace{0.4cm}

\textbf{A6.} Car P moves north at $40\ \text{km\,h}^{-1}$
and car Q moves south at $40\ \text{km\,h}^{-1}$. The
velocity of P relative to Q is:

\mcoption{A}{$0\ \text{km\,h}^{-1}$}
\mcoption{B}{$40\ \text{km\,h}^{-1}$ north}
\mcoption{C}{$80\ \text{km\,h}^{-1}$ north}
\mcoption{D}{$80\ \text{km\,h}^{-1}$ south}
\ansref{Ch2-A6}

\vspace{0.4cm}

\textbf{A7.} The resultant of two forces is maximum when
the angle between them is:

\mcoption{A}{$0°$}
\mcoption{B}{$90°$}
\mcoption{C}{$120°$}
\mcoption{D}{$180°$}
\ansref{Ch2-A7}

\vspace{0.4cm}

\textbf{A8.} A displacement is given as
$\vect{s} = 5\hat{i} - 12\hat{j}\ \text{m}$.
The magnitude of this displacement is:

\mcoption{A}{$7\ \text{m}$}
\mcoption{B}{$13\ \text{m}$}
\mcoption{C}{$17\ \text{m}$}
\mcoption{D}{$\sqrt{7}\ \text{m}$}
\ansref{Ch2-A8}

\vspace{0.4cm}

\textbf{A9.} Two equal forces $F$ act at right angles to
each other. Their resultant has magnitude:

\mcoption{A}{$F$}
\mcoption{B}{$F\sqrt{2}$}
\mcoption{C}{$2F$}
\mcoption{D}{$F/\sqrt{2}$}
\ansref{Ch2-A9}

\vspace{0.4cm}

\textbf{A10.} A particle is in equilibrium under three
concurrent forces. This means:

\mcoption{A}{All three forces must have the same magnitude}
\mcoption{B}{The resultant of any two forces equals the
third in magnitude but opposite in direction}
\mcoption{C}{All three forces must be parallel}
\mcoption{D}{The magnitude of the largest force must equal
the sum of the other two}
\ansref{Ch2-A10}

\vspace{0.4cm}

\textbf{A11.} The dot product $\vect{A}\cdot\vect{B}$
is zero when:

\mcoption{A}{$\vect{A}$ and $\vect{B}$ are parallel}
\mcoption{B}{$\vect{A}$ and $\vect{B}$ are perpendicular}
\mcoption{C}{$|\vect{A}| = |\vect{B}|$}
\mcoption{D}{$\vect{A}$ and $\vect{B}$ are anti-parallel}
\ansref{Ch2-A11}

\vspace{0.4cm}

\textbf{A12.} If $\vect{A} = 2\hat{i} + 3\hat{j}$ and
$\vect{B} = 4\hat{i} - \hat{j}$, then
$\vect{A}\cdot\vect{B}$ equals:

\mcoption{A}{$11$}
\mcoption{B}{$5$}
\mcoption{C}{$-5$}
\mcoption{D}{$8\hat{i} - 3\hat{j}$}
\ansref{Ch2-A12}

\vspace{0.4cm}

\textbf{A13.} The cross product $\vect{A}\times\vect{B}$
is a vector that is:

\mcoption{A}{Parallel to $\vect{A}$}
\mcoption{B}{Parallel to $\vect{B}$}
\mcoption{C}{Perpendicular to both $\vect{A}$ and $\vect{B}$}
\mcoption{D}{In the same plane as $\vect{A}$ and $\vect{B}$}
\ansref{Ch2-A13}

\vspace{0.4cm}

\textbf{A14.} $\hat{j}\times\hat{k}$ equals:

\mcoption{A}{$\hat{j}$}
\mcoption{B}{$-\hat{i}$}
\mcoption{C}{$\hat{i}$}
\mcoption{D}{$\hat{k}$}
\ansref{Ch2-A14}

\vspace{0.4cm}

\textbf{A15.} The unit vector in the direction of
$\vect{v} = 3\hat{i} + 4\hat{j}$ is:

\mcoption{A}{$3\hat{i} + 4\hat{j}$}
\mcoption{B}{$\dfrac{3}{5}\hat{i} + \dfrac{4}{5}\hat{j}$}
\mcoption{C}{$\dfrac{3}{7}\hat{i} + \dfrac{4}{7}\hat{j}$}
\mcoption{D}{$\hat{i} + \hat{j}$}
\ansref{Ch2-A15}

\vspace{0.6cm}

\subsection*{Section B: Theory and Calculations}

\textbf{B1.}
(a) Explain why two students cannot meaningfully add a
displacement vector and a force vector, even if both have
the same numerical value.\\[4pt]
(b) A hiker walks $4.0\ \text{km}$ due east, then
$3.0\ \text{km}$ due north, then $2.0\ \text{km}$ at
$45°$ north of east. Using the component method, find
the hiker's total displacement from the starting point
(magnitude and direction).
\hfill\ansref{Ch2-B1}

\vspace{0.4cm}

\textbf{B2.} A kite string makes an angle of $55°$ with
the ground. The tension in the string is $85\ \text{N}$.
A horizontal wind also pushes the kite with a force of
$30\ \text{N}$.\\[4pt]
(a) Find the horizontal and vertical components of the
string tension.\\
(b) Find the resultant of the string tension and the
wind force.\\
(c) In which direction does this resultant act?
\hfill\ansref{Ch2-B2}

\vspace{0.4cm}

\textbf{B3.} A boat that can travel at $5.0\ \text{m\,s}^{-1}$
in still water needs to cross a river $120\ \text{m}$ wide.
The river current flows at $3.0\ \text{m\,s}^{-1}$ due east.\\[4pt]
(a) If the boat points due north, what is its velocity
relative to the riverbank (magnitude and direction)?\\
(b) How far downstream does the boat land from its
intended destination?\\
(c) In what direction should the boat point to travel
directly north across the river, and how long does the
crossing take?
\hfill\ansref{Ch2-B3}

\vspace{0.4cm}

\textbf{B4.} Four forces act simultaneously on a metal
bracket:
\begin{itemize}[noitemsep]
  \item $\vect{F}_1 = 120\ \text{N}$ due east
  \item $\vect{F}_2 = 90\ \text{N}$ at $30°$ north of east
  \item $\vect{F}_3 = 150\ \text{N}$ due north
  \item $\vect{F}_4 = 80\ \text{N}$ at $45°$ south of west
\end{itemize}
(a) Resolve each force into $x$ (east) and $y$ (north)
components.\\
(b) Find the resultant force: magnitude and direction.\\
(c) What single fifth force would keep the bracket in
equilibrium?
\hfill\ansref{Ch2-B4}

\vspace{0.4cm}

\textbf{B5.} A $50\ \text{kg}$ traffic sign is suspended
by two cables attached to a horizontal beam. The left cable
makes an angle of $40°$ with the horizontal and the right
cable makes an angle of $55°$ with the horizontal.
(Take $g = 10\ \text{m\,s}^{-2}$.)\\[4pt]
(a) Draw a clear force diagram showing all forces on the
junction point.\\
(b) Write the equilibrium conditions ($\sum F_x = 0$ and
$\sum F_y = 0$).\\
(c) Solve for the tension in each cable.\\
(d) Which cable carries more tension? Explain why this
makes physical sense from the geometry of the problem.
\hfill\ansref{Ch2-B5}

\vspace{0.4cm}

\textbf{B6.} (Dot product application)\\
Two displacement vectors are
$\vect{d}_1 = 4\hat{i} + 3\hat{j}\ \text{m}$ and
$\vect{d}_2 = -2\hat{i} + 5\hat{j}\ \text{m}$.\\[4pt]
(a) Find $\vect{d}_1 \cdot \vect{d}_2$.\\
(b) Find the angle between $\vect{d}_1$ and $\vect{d}_2$.\\
(c) A constant force $\vect{F} = 10\hat{i} + 4\hat{j}\ \text{N}$
moves an object through displacement $\vect{d}_1$.
Find the work done by the force.\\
(d) Is there a direction in the $xy$-plane perpendicular
to $\vect{d}_1$? Write its unit vector.
\hfill\ansref{Ch2-B6}

\vspace{0.4cm}

\textbf{B7.} (Cross product application)\\
A rigid bar is pinned at the origin. A force
$\vect{F} = 8\hat{i} + 6\hat{j}\ \text{N}$
is applied at the point
$\vect{r} = 0.5\hat{i} - 0.2\hat{j}\ \text{m}$.\\[4pt]
(a) Calculate the torque $\vect{\tau} = \vect{r}\times\vect{F}$.\\
(b) State the magnitude of the torque and its direction
(clockwise or anticlockwise about the $z$-axis).\\
(c) Verify that $\vect{\tau}$ is perpendicular to $\vect{r}$
by computing $\vect{\tau}\cdot\vect{r}$.
\hfill\ansref{Ch2-B7}

\vspace{0.4cm}

\textbf{B8.} (Mixed vector operations)\\
Three vectors are given:
$\vect{A} = 2\hat{i} - \hat{j} + 2\hat{k}$,\;
$\vect{B} = \hat{i} + 3\hat{j} - \hat{k}$,\;
$\vect{C} = -\hat{i} + 2\hat{j} + 3\hat{k}$.\\[4pt]
(a) Find the unit vector $\hat{A}$ in the direction of
$\vect{A}$.\\
(b) Find $\vect{A}\cdot\vect{B}$ and hence the angle
between $\vect{A}$ and $\vect{B}$.\\
(c) Find $\vect{A}\times\vect{B}$.\\
(d) Verify that your answer to (c) is perpendicular to
both $\vect{A}$ and $\vect{B}$.\\
(e) Is $\vect{C}$ coplanar with $\vect{A}$ and $\vect{B}$?
(Hint: compute $\vect{C}\cdot(\vect{A}\times\vect{B})$;
zero means coplanar.)
\hfill\ansref{Ch2-B8}

\end{exercisebox}
